\begin{textsample}{POS Dim 3 – rs – Score 14.00 – rs/rs\_inf\_16.txt}  \label{ex:f3_pos_250}
11.4.1 Direitos de Aprendizagem e Desenvolvimento no Campo de Experiências \textbf{Escuta}, Fala, Pensamento e Imaginação CONVIVER com \textbf{crianças} e \textbf{adultos} em situações comunicativas cotidianas, constituindo modos de pensar, imaginar, sentir, narrar, dialogar e conhecer. \textbf{BRINCAR} com parlendas, trava-línguas, adivinhas, memória, \textbf{rodas}, \textbf{brincadeiras} cantadas, jogos e textos de imagens, escritos e outros, ampliando o repertório das manifestações culturais da tradição local e de outras culturas, enriquecendo sua linguagem oral, corporal, musical, \textbf{dramática}, escrita, dentre outras.
PARTICIPAR de \textbf{rodas} de conversa, de relatos de experiências, de contação e leitura de histórias e poesias, de construção de narrativas, da elaboração, descrição e representação de papéis no faz de conta, da exploração de materiais impressos e de variedades linguísticas, construindo diversas formas de organizar o pensamento.
EXPLORAR \textbf{gestos}, expressões, \textbf{sons} da língua, rimas, imagens, textos escritos, além dos sentidos das palavras, nas poesias, parlendas, canções e nos enredos de histórias, apropriando -se desses elementos para criar novas falas, enredos, histórias e escritas convencionais ou não.
EXPRESSAR sentimentos, ideias, \textbf{percepções}, \textbf{desejos}, necessidades, pontos de vista, informações, dúvidas e descobertas, utilizando múltiplas linguagens, considerando o que é comunicado pelos \textbf{colegas} e \textbf{adultos}.
CONHECER -SE e reconhecer suas preferências por pessoas, \textbf{brincadeiras}, lugares, histórias, autores, gêneros linguísticos, e seu interesse em produzir com a linguagem verbal. 11.4.2 Objetivos de Aprendizagem e Desenvolvimento Os Objetivos de aprendizagem e desenvolvimento para os \textbf{bebês}, as \textbf{crianças} bem \textbf{pequenas} e as \textbf{crianças} \textbf{pequenas} assim são descritos na BNCC e no Referencial Curricular Gaúcho :

% matched lemmas: adulto, bebê, brincadeira, brincar, colega, criança, desejo, dramático, escutar, gesto, pequeno, percepção, roda, som
\end{textsample}
